\documentclass[../main.tex]{subfiles}

\begin{document}
	
	\chapter*{Introduzione}
	\addcontentsline{toc}{chapter}{Introduzione} \markboth{INTRODUZIONE}{}
	
	Comprendere e prevedere il comportamento della viscosità dei liquidi sottoraffreddati rappresenta una sfida importante sia per la fisica della materia condensata sia per l'industria dei materiali. Man mano che un liquido viene raffreddato verso la temperatura di transizione vetrosa, la sua viscosità macroscopica aumenta di oltre dodici ordini di grandezza. Mappare sperimentalmente questo andamento lungo l'intero intervallo termico per nuove composizioni chimiche è un processo oneroso e tecnicamente complesso. Per tale motivo, negli ultimi anni, lo sviluppo di modelli predittivi \textit{data-driven} ha assunto un ruolo centrale, permettendo di stimare le proprietà reologiche dei vetri direttamente a partire dalla loro composizione nominale.
	
	Nel panorama del \textit{machine learining} applicato alla scienza dei vetri, le reti neurali hanno dimostrato capacità predittive eccezionali. Il punto di riferimento fondamentale per questo lavoro di tesi è il modello ViscNet, sviluppato da Cassar~\cite{cassarViscNetNeuralNetwork2021}. ViscNet adotta un interessante approccio \textit{gray-box}: anziché prevedere direttamente la viscosità, la rete neurale stima i tre parametri termodinamici dell'equazione analitica di MYEGA ($T_g$, $m$, $\log_{10}\eta_\infty$). Questo vincolo garantisce coerenza fisica alle previsioni.
	
	Tuttavia, il ricorso ad architetture di \textit{deep learning} porta con sé un limite significativo: la perdita di interpretabilità. A causa delle stratificazioni e delle attivazioni non lineari interne, modelli come ViscNet operano di fatto come "scatole nere". Pur fornendo previsioni esatte, non permettono di comprendere in modo trasparente quali specifiche proprietà chimico-fisiche degli elementi costituenti guidino l'arresto dinamico del liquido. Metodi di interpretabilità \textit{post-hoc} (come l'analisi SHAP) offrono solo letture locali ed euristiche, lasciando irrisolto il problema della comprensione globale del fenomeno.
	
	Da questo limite nasce la motivazione principale di questo studio. L'obiettivo è verificare se sia possibile sostituire l'architettura neurale con modelli di regressione lineare interpretabili, senza sacrificare l'accuratezza predittiva. In un modello lineare ideale, infatti, ciascun coefficiente quantifica in modo diretto, globale e inequivocabile l'impatto di una singola variabile sulla grandezza target.
	
	Per garantire un confronto rigoroso, in questo lavoro si è assunto il modello ViscNet come \textit{benchmark} assoluto. Si è riprodotta fedelmente l'intera \textit{pipeline} di partenza utilizzata da Cassar: dal dataset originale (deduplicato dal database SciGlass), alla divisione tra dati di \textit{training} e di \textit{test}, fino all'estrazione dello spazio ad alta dimensionalità composto da 560 descrittori chimico-fisici. 
	
	Su questa base, sono state addestrate e ottimizzate diverse varianti di regressione lineare multivariata (Ridge, Elastic Net, VIF+OLS e regressioni su componenti principali come PCR e SPCR) per mitigare il forte problema della multicollinearità intrinseca ai dati chimici. Attraverso l'analisi comparativa delle prestazioni e lo studio dei pesi assegnati dai vari algoritmi, la tesi si propone non solo di testare il limite predittivo dei modelli lineari in questo dominio, ma anche di indagare se l'interpretabilità matematica si traduca effettivamente in una chiara interpretabilità fisica del fenomeno analizzato.
	
\end{document}